\section{Introduction}
Canonical forms serve as a cornerstone for systematic feedback
design in finite dimensions. Their counterparts for \glspl{dps}
have been studied since the late 1970s. 
For linear hyperbolic \glspl{pde}, \cite{Russel1978}
connected controllability and spectral assignment to
\glspl{ccf} and neutral delay equations.  
More recent works linked the aforementioned delay interpretations to flatness
based parameterisations \cite{FMRR95cdc,FliMou99} -- see also \cite{CookeKrumme1968jmaa} for early ideas in this direction -- and, even more recently, established the connection between such parameterisations and the \gls{ccf} for more general systems.
In particular, \cite{Woittennek2012-ifacwc} introduced a 
\gls{hccf} for linear hyperbolic SISO systems comprising general interconnections of hyperbolic \gls{pde} and \gls{ode} (see also \cite{Woittennek2011ifac}). This \gls{hccf}
generalises the familiar chain-of-integrators
structure by prepending a transport equation to represent the delays
of the system imposed by the hyperbolic \glspl{pde}.
Therein, the canonical coordinates are introduced on the basis of a \gls{fde} 
arising from a flatness-based parametrisation of the system's input while the required state transforms are computed using the flatness-based parametrisation of the system's original state variables.
Aside from the above-mentioned state-space approaches, in  \cite{Woittennek2010-siam,Gehring2013}  flatness-based parametrisations also served as means for purely algebraic analysis and control design. These methods employ and partially generalise results from
\cite{Brethe1996} and \cite{GluesingLuerssen1997}.

As many relevant technological applications posses multiple in- and outputs,
several control approaches are available, e.g.~ \cite{Knueppel2014} using
flatness-based methods or \cite{Gabriel2024} employing backstepping designs.
However, regarding controller forms a direct extension of the existing work to
the \acs{mimo} case is difficult, as such systems can exihibit
different delays depending on the considered input which together with the
bidirectionally coupled ODE result an ill-posed \gls{fde} system,
thus negating the approach from the \acs{siso} case.

Therefore, the goal of this contribution is to introduce a \gls{hcf} for a class
of linear hyperbolic \acs{mimo} systems with dynamic boundary conditions and to
develop a delay reduction algorithm such that the \gls{fde} system is well
posed and the controller form can be computed.
To do so, we denote the input-output relation using quasipolynomials whose use
in the control of hyperbolic systems was already reviewed in \cite{Loiseau1998}
as well as \cite{Benarab2022} and restate the requirements to introduce a \gls{hcf}
as column-reduced properties on the resulting quasipolynomial matrices.
Next, in contrast to the established methods and due to the incommensurate
delays in the considered systems as well as the interplay with the boundary
system, we propose a modified quasipolynomial division algorithm to compute
the required transformation to establish said properties. 
% hyperbolic systems with boundaries governed by \acsp{ode} for the \acs{siso}
% case which this contribution extends to the \acs{mimo} case.


% Backstepping provides an alternative, transformation-based route that leverages
% strict-feedback structure and recursive stabilization. Importantly for the
% present MIMO focus, such designs are not limited to scalar inputs: extensions
% cover multi-input–multi-output ODE–PDE–ODE systems with general
% heterodirectional hyperbolic subsystems and additional boundary ODEs, as well
% as underactuated interconnections and parabolic variants. However, while
% state-feedback stabilization is systematic, observer and output-feedback
% designs often hinge on approach-specific structures and measurement placements
% rather than intrinsic normal forms, with some cases being dual to state
% feedback and others substantially more involved.


% From the onset of state space based approaches, 
% controller (canonical) forms have been a vital asset
% for systematic controller designs.
% As the controller forms for linear \acs{siso} and \acs{mimo} \acs{ode} systems,
% cf.~\cite{Kalman1963,Luenberger1967} as well for \acs{siso} hyperbolic
% \acs{pde} systems \cite{Russel1978} are long established and understood,
% only comparatively recent works 
% \cite{Woittennek2012-ifacwc,Woittennek2013-cpde,Gehring2023} 
% discuss
% hyperbolic systems with boundaries governed by \acsp{ode} for the \acs{siso}
% case which this contribution extends to the \acs{mimo} case.

Following some preliminaries, where we briefly restate established controller 
forms and demonstrate the problems that arise if a hyperbolic controller form
is to be introduced even for the most simple \acs{mimo} system,
a general \acs{mimo} \gls{hcf} and several subforms are given in \Cref{sec:mimo_hcf}.
The main part of the contribution in \Cref{sec:computation} then discusses
how a \gls{hcf} can be obtained for a specific class of systems 
Finally, \Cref{sec:example} tracks the procedure for an example system.


\section{Preliminaries}

% \subsubsection*{Notation}
\subsection{Notation}

The variables $z$ and $t$ denote space and time.
The first and second derivative of $x(t)$ is denoted
by $\dot x(t)$ and $\ddot x(t)$, while the $n$-th 
derivative is given by $x^{(n)}(t)$
% Regarding vectors, the operation
% $\bm x^{[\kron]}(t)$ is used  to substitute each 
% $x_i(t)$ in $\bm x(t)$ by $x_i^{(\kron_i)}(t)$.
% For unknowns in several variables 
% Furthermore, we write e.g.~for 
and the partial derivative of $x(z,t)$ w.r.t.~$z$.
by $\partial_z x(z,t)$.


To treat \glspl{fde} in an algebraic setting, we identify the
operators therein with their Fourier-Laplace transforms 
and treat the application of the transformed operators
as convolutions.
Thus, we write $\mathcal{L}\{\diffsym\} = \lapsym$ for
differentials, while  
the shifts
$(\shiftsym^\tau x)(t) = x(t+\tau)$ 
with $\tau \in \reals$
are denoted as
$\mathcal{L}\{\shiftsym^\tau\}(\lapsym) = \exp(\lapsym\tau)$.
% In that context,
% $\mathcal{L}\{\flatsym(t)\} = \lapflatsym(s)$.
% the Laplace transform of e.g.~$\flatsym(t)$ is denoted by $\lapflatsym(s)$.

The shorthand $\leb([a,b])$ is used for the space of square integrable
functions on $[a,b]$, while $\sob[n]([a,b])$ stands for the space of functions whose $n$-th
derivative belongs to $\leb([a,b])$
and $\entirefuncs$
% as well as $\lapentirefuncs$ 
for Schwartz-Distributions with compact support.
% and their Laplace transforms, respectively
% occurring as the Laplace transformations of shift operators. 
With $\tvar_f([a,b])$ denoting the total variation of a function $f$ on the
interval $[a,b]$, $\boundvar([a,b])$ contains functions of bounded
variation while members of its subset $\boundvarleft([a,b])$ 
satisfy
$\lim_{\epsilon \to 0} \tvar_f([b,b+\epsilon)) = 0$.
% $\boundvarleft([a,b]) = \{
% 	f \in \boundvar([a,b]) 
% 	\vert 
% 	\lim_{\epsilon \to 0} \tvar_f([a,a+\epsilon)) = 0
% \}$

Furthermore, $\realpolys$ is the space of polynomials in $\diffsym$ with coefficients in $\reals$
% , $\realgpolys = \mathbb{R}\left[\ggen\right]$ contains generalised polynomials with real-valued exponents 
and $\realrationals$ consists of rational expressions in $\diffsym$. 
Regarding a polynomial matrix $X(\diffsym) \in \reals^{n\times n}[\diffsym]$, we denote the \gls{lccm} w.r.t.~$\diffsym$ by $\lccm_{\diffsym}{\mat X(\diffsym)}$ and say that $\mat X(\diffsym)$ is column reduced w.r.t.~$\diffsym$ if $\lccm_{\diffsym}{\mat X(\diffsym)}$ has full generic rank.
Finally, we call $\mat X(\diffsym)$ unimodular if its determinant is a nonzero constant (and therefore independent of $\diffsym$).


